\section{Uniform bounds for the link diagrams}\label{bd:section}

This section supplies the uniform bounds used in the Sarkisov argument.
Theorem~\ref{bd:exclusions} converts bounded varieties and marked covers
into degree bounds for K3 tests. We then construct a central model at
each Sarkisov wall, preserving a fixed discrepancy bound, and recover
the link as contractions of that model. Theorem~\ref{bd:diagrams}
bounds these diagrams on geometric generic fibres, including the maps
and marked bad loci that control the branch locations of later covers.
Theorem~\ref{bd:terminalization} separately bounds models extracting
only low-discrepancy valuations from a bounded log-Fano pair; this will
control a relative minimal model without bounding its resolving start.
Finally, Corollary~\ref{bd:choice-order} first fixes
one ramification coefficient from the ordinary fourfold diagrams;
only then do we bound the auxiliary log diagrams and choose the final
K3 degree threshold.

\subsection{Elementary exclusions for K3 fields}

In this subsection a \emph{degree-$d$ K3 surface} means a smooth projective
K3 surface $S$ with $\operatorname{Pic}(S)=\Z L$, where $L$ is ample and
$L^2=d$.  All the bounds below are independent of the choice of $S$, and
therefore also of any further restriction on its transcendental Hodge
structure or automorphism group.

\begin{lemma}[Pencils on a K3 surface]\label{bd:pencil}
Let $S$ be a degree-$d$ K3 surface.  If a dominant rational map from $S$
to a smooth projective curve has connected general fibre, that curve is
$\mathbb P^1$.  If $g$ is the genus of the smooth general fibre on a
resolution of the map, then
\[
                         2g-2\geq\sqrt d.
\]
For an arbitrary dominant rational map to a curve the same inequality
holds for every geometric general fibre component, after normalization.
\end{lemma}

\begin{proof}
Resolve the map by point blow-ups $\pi:Y\to S$.  Since $h^1(Y,\mathcal
O_Y)=0$, pullback of differentials shows that the base of its Stein
factorization has genus zero.  Thus, after Stein factorization, we have
a morphism $f:Y\to\mathbb P^1$ with smooth connected general fibre $F$.
The images on $S$ of these fibres form a pencil without fixed components,
of class $nL$ for some integer $n\geq1$.  In the total-transform basis of
the successive exceptional curves write
\[
 F=\pi^*(nL)-\sum_i m_iE_i^*,\qquad
 K_Y=\sum_iE_i^*,\qquad m_i\geq0.
\]
The exceptional basis is orthogonal, with $(E_i^*)^2=-1$.  Consequently
adjunction and $F^2=0$ give
\[
       \sum_i m_i^2=n^2d,\qquad 2g-2=K_Y\cdot F=\sum_i m_i.
\]
It follows that $2g-2\geq(\sum_i m_i^2)^{1/2}=n\sqrt d$.
Stein factorization accounts for the geometric components in the last
assertion.
\end{proof}

We use the following precise boundedness conventions.  A collection of
projective varieties is \emph{bounded} if its members occur as geometric
fibres of finitely many projective morphisms of schemes of finite type.
It is \emph{birationally bounded} if each member is birational to such a
fibre.  For pairs $(Y,D)$, boundedness includes the reduced proper closed
subset $D\subsetneq Y$: it must be supplied by a closed subset of the
same family.  A finite cover of $(Y,D)$ means the normalization of $Y$
in a finite field extension, required to be \emph{\'etale over $Y\setminus
D$}.  Thus the boundary hypothesis controls branch locations, not just
the degree of the extension.

Here and below one can choose a bounded collection of smooth
projective birational models of the irreducible components of a bounded
collection.  To justify this
operation, work on an irreducible parameter space, separate its geometric
generic components by a finite extension of the function field, resolve
them, and spread the resolutions.  After shrinking, the spread models
are smooth and the maps are birational on the appropriate fibre
components.  Apply the same construction to the remaining closed subset.
Noetherian induction terminates this procedure.  The same argument works
with a marked closed subset, using embedded resolution and including the
exceptional locus in the new boundary.  It also proves that irreducible
components of fibres of a fixed projective family, and smooth projective
models of those components, are birationally bounded.  These operations
allow finite base changes of parameter spaces; they impose no bound on a
particular map from one member to another.

\begin{lemma}[Bounded threefolds and their K3 quotients]
\label{bd:bounded-mrc}
For a birationally bounded collection of varieties of dimension at most
three, there is an integer $b$ such that any degree-$d$ K3 surface
birational to the MRC base of a member satisfies $d\leq b$.
\end{lemma}

\begin{proof}
First consider a bounded smooth projective surface $Y$ birational to
$S$.  Its minimal-model morphism $\mu:Y\to S$ is a succession of point
blow-downs.  Let $A$ be a very ample divisor from the bounded family.
In the total-transform exceptional basis,
\[
 A=\mu^*B-\sum_i a_iE_i^*,\qquad
 K_Y=\sum_iE_i^*,\qquad a_i=A\cdot E_i^*>0.
\]
The pushforward $B$ is ample: at each blow-down its square increases and
its intersection with every curve is positive, so the surface
Nakai--Moishezon criterion applies \cite[Theorem~1]{Cartier66}.  It is an integral multiple of $L$.
Hence
\begin{equation}\label{bd:surface-degree}
 d\leq B^2=A^2+\sum_i a_i^2
       \leq A^2+(A\cdot K_Y)^2.
\end{equation}
The right-hand side takes only finitely many values in a bounded smooth
projective family after a finite stratification.  This proves the
assertion in dimension two; in smaller dimensions a K3 MRC base is
impossible.

It remains to bound the birational types of the K3 bases of bounded
smooth projective threefolds.  Put these threefolds into finitely many
smooth projective families $\mathcal X\to T$, with relatively very ample
divisor $H$. We use the fixed-polynomial Hilbert space and its open
graph locus \cite[Tags~0DPH, 0D1A and~0D1B]{StacksProject057}, together with
semicontinuity and base change \cite[Tags~0BDN and~0B91]{StacksProject057}.
For $e\geq1$, the free locus in the relative scheme of
degree-$e$ maps $\mathbb P^1\to\mathcal X/T$ is open.  Its morphism to
$T$ is smooth: for a free map $a$ the obstruction space
$H^1(\mathbb P^1,a^*T_{\mathcal X_t})$ vanishes.  Its image is therefore
open.  If $U_D$ is the union of these images for $1\leq e\leq D$, then
the $U_D$ are ascending open subsets of the noetherian space $T$.
Their union is precisely the locus of uniruled geometric fibres, since
in characteristic zero a smooth projective variety is uniruled if and
only if it has a nonconstant free rational curve.  The union is
quasi-compact, so $U_D$ equals that union for some $D$.
These deformation statements follow, for example, from
\cite[Remark~9, Proposition~10 and Theorem~15]{AraujoKollar02}; this argument proves
the uniformity rather than assuming that the uniruled fibres themselves
form a separately specified family.

Suppose now that the MRC quotient of $X=\mathcal X_t$ is a surface
birational to $S$.  Use its almost-holomorphic presentation
$r:U\to S^\circ$ with proper smooth rationally connected fibres;
see \cite[Chapter~IV, Section~5]{Kollar96}.
Its fibres are smooth copies of $\mathbb P^1$.  A degree-at-most-$D$
free map on $X$ deforms in a covering family.  A very general member
meets a very general fibre of $r$ and is contained in it, by the MRC
property.  Since that fibre is a complete integral curve, the image of
the map is the whole fibre.  The $H$-degree $e$ of the fibre is thus at
most $D$.

For such a fibre $F$, smoothness of $r$ gives
$N_{F/X}\simeq\mathcal O_{\mathbb P^1}^{\oplus2}$.  Its Hilbert scheme
is therefore smooth of dimension two at $[F]$.  The family of fibres
defines an injective map from $S^\circ$ to this Hilbert scheme, with
isomorphic tangent spaces; its closure is an irreducible component of
$\operatorname{Hilb}^{en+1}(X)$ birational to $S$.  The components so
obtained are components of fibres of the finite-type projective scheme
\[
              \coprod_{e=1}^{D}
                 \operatorname{Hilb}^{en+1}(\mathcal X/T).
\]
They have bounded smooth projective birational models by the preceding
bounded-family reduction.  Apply \eqref{bd:surface-degree} to these
surfaces.  This proves a uniform bound independent of the threefold and
of its MRC quotient map.
\end{proof}

\begin{theorem}[Bounded data exclude large K3 degrees]
\label{bd:exclusions}
Fix the finite-type families occurring in each of the following
statements, and fix an integer $N\geq1$.  There is an integer $b$,
depending only on these families and on $N$, with the following
properties.
\begin{enumerate}[label=\textup{(\roman*)}]
\item No member of a specified birationally bounded collection of
varieties of dimension at most three has a degree-$d$ K3 MRC base with
$d>b$.
\item Let $(Y,D)$ range over a bounded collection of projective curve
or surface pairs.  The smooth projective curves covering the curve
members with degree at most $N$, \emph{\'etale off $D$}, have bounded
genus.  No surface field obtained by such a cover of a surface member is
the field of a degree-$d$ K3 surface for $d>b$.
\item Let $T\dashrightarrow C$ be a dominant map from an integral
complex surface to a smooth projective curve.  Suppose each normalized
geometric generic fibre component is a degree-at-most-$N$ cover of a
member of a fixed bounded collection of curve pairs, \emph{\'etale off
the marked boundary}.  Then $T$ is not birational to a degree-$d$ K3
surface for $d>b$.
\item Let $T\dashrightarrow C$ be a dominant map from an integral
complex threefold to a smooth projective curve.  Suppose the geometric
generic fibre components are uniruled and belong to a fixed
birationally bounded collection of surfaces.  Then the MRC base of
$T$ is not a degree-$d$ K3 surface for $d>b$.
\end{enumerate}
The bounded geometric data in \textup{(iii)} and \textup{(iv)} are
understood after algebraic closure of $\C(C)$; neither the total spaces
$T$ nor the maps to $C$ are required to be bounded.  The assertions
apply to each field factor of a finite algebra separately.
\end{theorem}

\begin{proof}
Part~\textup{(i)} is Lemma~\ref{bd:bounded-mrc}.
For \textup{(ii)}, pass to bounded smooth projective models of the pairs,
including the exceptional locus in the boundary.  The normalized
base-changed covers are still \emph{\'etale} on the complementary open.
For a smooth curve $A$ of genus $g_A$ with at most $r$ marked points,
and a connected degree-$e$ cover $A'\to A$ unramified elsewhere,
Riemann--Hurwitz gives
\begin{equation}\label{bd:hurwitz}
 2g(A')-2\leq e(2g_A-2)+(e-1)r.
\end{equation}
Indeed, over any one branch point the ramification contribution is
$e$ minus the number of points above it, and is at most $e-1$.
This proves the curve assertion uniformly for $e\leq N$.

For a bounded smooth surface pair choose a pencil of hyperplane sections
for a relatively very ample divisor $H$.  A general section $A$ is smooth,
has uniformly bounded genus by adjunction, and meets the divisorial part
of $D$ in uniformly boundedly many points.  Choose it to avoid the
zero-dimensional part of $D$.  For any one cover, general members can
also avoid the images of the singular points of its normalization;
normal surfaces have only finitely many singular points.  Each
normalized component over $A$ has degree at most $N$ and is ramified
only over $A\cap D$, so \eqref{bd:hurwitz} gives a uniform genus bound
$G$.  Resolve the pulled-back pencil on a smooth model of the covering
surface and take its Stein factorization.  Its smooth connected general
fibres are precisely these normalized components.  If the surface field
were that of $S$, Lemma~\ref{bd:pencil} would give
$\sqrt d\leq2G-2$.  This proves the surface assertion.  The pencil can
depend on the individual cover; only its genus and boundary-intersection
bounds need to be uniform.

The same curve estimate over $\overline{\C(C)}$ bounds the genera of
the geometric generic components in \textup{(iii)}.  If $T$ were
birational to $S$, the given map to $C$, after resolution and Stein
factorization, would give a pencil on $S$.  Lemma~\ref{bd:pencil}
again bounds $d$.  No genus bound for the parameter curve $C$ is needed.

For \textup{(iv)}, suppose the MRC base of $T$ is birational to $S$,
and write $r:T\dashrightarrow S$ and $f:T\dashrightarrow C$.
Put $F=\C(C)$ and choose an algebraic closure $\overline F$.  Resolve
the maps when necessary.  If $f$ were nonconstant on a general fibre
of $r$, the map $(r,f):T\dashrightarrow S\times C$ would be dominant
and generically finite.  A component of a geometric generic surface
fibre of $f$ would then dominate $S_{\overline F}$ generically finitely.  Such a
component is uniruled, whereas a generically finite dominant image of
an uniruled surface is uniruled.  This contradicts the fact that a K3
surface is not uniruled in characteristic zero.

Consequently $f=h\circ r$ for a dominant rational map $h:S\dashrightarrow
C$.  This is descent along the geometrically integral general fibres
of the MRC quotient, or equivalently the relative algebraic closedness
of $\C(S)$ in $\C(T)$.  Resolve $h$ and take its Stein factorization.
Over $\overline F$, let $\Gamma$ be a smooth geometric generic fibre
of the resulting connected-fibre pencil, of genus $g$.  A geometric
generic component $P$ of $f$ dominates the corresponding $\Gamma$.
For a smooth projective model $\widetilde P$ over $\overline F$ this implies
\[
 g\leq h^0(\widetilde P,\Omega^1_{\widetilde P})
       =h^1(\widetilde P,\mathcal O_{\widetilde P}).
\]
To see the inequality, resolve the rational map to $\Gamma$ and pull
back regular one-forms; pullback is injective for a dominant map in
characteristic zero, and irregularity is a birational invariant of
smooth projective surfaces.  The irregularities of $\widetilde P$ are
uniformly bounded by bounded smooth projective birational models and
upper semicontinuity.  Lemma~\ref{bd:pencil} therefore bounds $d$.

All genus, intersection, and coherent-cohomology bounds used here are
preserved by extension of algebraically closed fields.  The same proofs
therefore apply to the geometric generic data specified in
\textup{(iii)} and \textup{(iv)}.  Taking the maximum of the finitely many
bounds proves all the assertions.
\end{proof}

\subsection{Central models for the two Sarkisov programs}

The boundedness problem starts with a wall in a Sarkisov factorization.
We first construct a central model that retains the prime valuations of
the link and a fixed lower bound for log discrepancies. Its geometric
generic fibre will supply the log-Fano pair to which boundedness is
applied. The next subsection will then recover the contraction maps and
their marked loci from these pairs.

For a pair $(Y,B)$ we write $a(F,Y,B)$ for \emph{log discrepancy}; a
prime divisor on $Y$ of coefficient $b$ has log discrepancy $1-b$.
A pair is $\epsilon$-lc if every log discrepancy is at least
$\epsilon$.  Terminality concerns exceptional divisors and means
$a(F,Y,B)>1$ for such divisors.  In particular, terminality of the
underlying variety is used for the ordinary Sarkisov program, whereas
only a fixed positive lower bound for log discrepancies is needed for
the auxiliary log program.

A variety is of \emph{Fano type} if it admits an effective rational
boundary making it klt log Fano.  We use the relative version over a
base as well.  By finite generation, a $\Q$-factorial variety of Fano
type is a Mori dream space: its movable cone has finitely many
semiample chambers, and an MMP for any divisor terminates
\cite[Corollary~1.3.2]{BCHM10}.  The chamber description includes
the contraction associated with each face
\cite[Proposition~1.11]{HuKeel00}.
The contraction and relative basepoint-free statements we use are
\cite[Theorem~4.5.2 and Corollary~6.9.4]{FujinoFoundationsDraft057}.
For the sign and uniqueness of exceptional differences we use the
negativity lemma \cite[Lemma~11.60]{KollarFamiliesDraft057}.

A \emph{contraction morphism} is a projective surjective morphism
$f:Y\to Z$ of normal varieties with $f_*\mathcal O_Y=\mathcal O_Z$.
A \emph{birational contraction} is a birational map whose inverse
contracts no divisor.
A \emph{small modification} is a birational map that is an isomorphism
in codimension one. A \emph{small $\Q$-factorialization} is a projective
small birational morphism with $\Q$-factorial source.
For a $\Q$-factorial Fano-type source $Y$, a
\emph{rational contraction} means a composite
$Y\dashrightarrow Y'\to Z$, where the first map is a small
$\Q$-factorial modification and the second is a contraction morphism.
The target $Z$ need only be normal and projective. We use the same
definition over an indicated base, with both maps over that base.

For an integral Weil divisor $D$ on a normal projective variety $Y$,
write
\[
 R(Y,D)=\bigoplus_{m\geq0}H^0(Y,\mathcal O_Y(mD)),
\]
where $\mathcal O_Y(mD)$ is the divisorial reflexive sheaf. For a
rational divisor we use a sufficiently divisible integral multiple;
further Veronese subrings do not change $\operatorname{Proj}$. In
the finitely generated cases used below, the induced rational map
$Y\dashrightarrow\operatorname{Proj}R(Y,D)$ is its \emph{ample-model map}.

We use the ample-model formulation of the Sarkisov theorem
\cite[Theorem~1.3, Theorem~3.3, Theorem~3.7, and Lemma~4.1]{HM13}.
Start with two Mori fibre-space outcomes of one klt $K_W+\Phi$
program on a smooth projective variety $W$. The theorem connects
them using a finite-dimensional family of klt boundaries $\Theta$
with $\Theta-\Phi$ ample on $W$; this family and its ample models
are the \emph{geography} used below.
All paths whose diagrams we bound below are chosen from this geography.

At a linking wall, let $N_1,N_2$ be the $\Q$-factorial middle models
on the two sides. The local diagram supplied by the two-ray construction is
\[
\begin{tikzcd}[row sep=small,column sep=large]
 N_1 \arrow[rr,dashed,"\text{small over }Q"] \arrow[d,"e_1"]
   && N_2 \arrow[d,"e_2"] \\
 V_1 \arrow[d,"\pi_1"] && V_2 \arrow[d,"\pi_2"] \\
 B_1 \arrow[dr,"q_1"'] && B_2 \arrow[dl,"q_2"] \\
 & Q &
\end{tikzcd}
\]
Each $N_i\to Q$ has relative Picard rank two. The arrow $e_i$ is
either the identity or an elementary divisorial contraction;
$\pi_i$ is the end Mori fibre contraction, of relative rank one.
Additivity of the relative ranks in these morphism towers gives
\begin{equation}\label{bd:rank-tower}
 2=\rho(N_i/Q)=\rho(N_i/V_i)+1+\rho(B_i/Q).
\end{equation}
These are ranks over the original common base $Q$, before any
extension of its function field. The arrows $q_i$ may be small
contractions. All maps $W\dashrightarrow N_i$ are birational
contractions, and the transforms of $K_W+\Theta$ on the $N_i$ are
semiample and pulled back from $Q$. Thus the middle small models
share their prime valuations; the divisorial exits $e_i$ can remove
valuations.

The uniformity we need is a deduction from this construction,
not an assumption that arbitrary Fano-type varieties are bounded.

A relative negative program over a projective target is also a
negative program in the absolute category used here.  Indeed its
cone of contracted curves is the face cut out by the pullback of
an ample divisor on that target; a negative extremal ray of this
face is extremal in the full cone.  The contractions and flips
remain projective.  Thus a program run first over one projective
target and then over another, for example over $T$ and then over
$Q$, is an allowed common-start program.  Its final relative
rank-one fibre contraction is a Mori fibre space in this category.

\begin{lemma}[A bounded-discrepancy central model]
\label{bd:central}
Consider a wall in this geography on a smooth projective start
$(W,\Phi)$.  Assume that $(W,\Phi)$ is $\epsilon$-lc.  Then the
relative small models in the link have a common set of prime
valuations and admit a small modification $N^{\mathrm c}/Q$
such that
\[
  -(K_{N^{\mathrm c}}+\Phi_{N^{\mathrm c}})
       \quad\text{is nef and big over }Q.
\]
Its log-anticanonical ample-model map
$N^{\mathrm c}\to Z^{\mathrm c}$ is small.  Both
$N^{\mathrm c}$ and $Z^{\mathrm c}$ are $\epsilon$-lc with
their transformed boundaries.  All the link vertices and
contraction bases are rational contractions of a small
$\Q$-factorialization of $Z^{\mathrm c}$.  If $\Phi=0$, all
the small models can be taken terminal.
\end{lemma}

\begin{proof}
Take $N=N_1$ in the displayed wall diagram, let
$A=\Theta-\Phi$ be the ample difference on $W$, and write
$A_N,\Theta_N,\Phi_N$ for transforms. We first transfer the
discrepancy bound from $W$ to $N$. We then make the original
log-anticanonical class nef and pass to its ample model. That last
map must be small: a divisorial contraction would lose one of the
prime valuations we need to retain.
On a graph resolution
with projections $p$ to $W$ and $q$ to $N$,
\begin{align}
 p^*(K_W+\Theta)-q^*(K_N+\Theta_N)&=E_\Theta\geq0,
       \label{bd:wall-comparison}\\
 p^*A-q^*A_N&\leq0.\label{bd:ample-comparison}
\end{align}
For this comparison at the wall, approach $\Theta$ by interior parameters
$\Theta_i$ in the chamber defining $N$.
By \cite[Theorem~3.3(3)]{HM13}, $W\dashrightarrow N$ is a log terminal
model for $K_W+\Theta_i$, so the exceptional difference on the same graph
resolution is effective. Its finitely many coefficients depend linearly
on the boundary. Taking their limits gives $E_\Theta\geq0$;
\cite[Theorem~3.3(2)]{HM13} supplies the contraction $N\to Q$.
Thus the first difference is exceptional and nonnegative at the wall.  The second is also $q$-exceptional,
since the map does not extract.  On a $q$-contracted curve its
degree equals that of $p^*A$, so it is $q$-nef; the negativity
lemma proves \eqref{bd:ample-comparison}.  Subtracting gives
\begin{equation}\label{bd:original-comparison}
 p^*(K_W+\Phi)-q^*(K_N+\Phi_N)\geq0.
\end{equation}
It follows valuation by valuation that the original
$\epsilon$-lc bound is retained on $N$.

There is a useful extension of this argument.  The divisor
$K_N+\Theta_N$ is pulled back from $Q$.  Any small
$\Q$-factorial modification $N'\to Q$ has the same
codimension-one divisor and the same pullback relation.
Thus $N\dashrightarrow N'$ is crepant for $K_N+\Theta_N$.
The comparison \eqref{bd:wall-comparison} remains nonnegative
for $W\dashrightarrow N'$, and the ample subtraction proves
\eqref{bd:original-comparison} for $N'$ as well.  In particular,
the bound holds on any required relative small chamber; there
is no need to locate all those chambers in one chosen
two-dimensional slice.

The pair $(N,\Theta_N)$ is klt, its boundary is big over $Q$,
and its log-canonical divisor is trivial over $Q$.  Writing
that boundary, up to relative rational linear equivalence, as
an ample divisor plus an effective divisor, a sufficiently
small perturbation makes a klt log-Fano boundary.  Therefore
$N/Q$ is of Fano type.  Furthermore,
\[
                 -(K_N+\Phi_N)\equiv_Q A_N.
\]
The transform of an ample class by a birational contraction is
in the interior of the movable cone: choose divisor classes
on $W$ mapping onto a basis on $N$, perturb $A$ slightly in
both directions in this basis, and push their ample systems
forward.  No prime divisor is fixed, because none is extracted.
The same argument works after passing to the relative numerical
space over $Q$.

Finite generation now makes this movable class nef on a small
model $N^{\mathrm c}/Q$.  Only small steps are needed: a
negative divisorial contraction would force its exceptional
divisor to be fixed in the corresponding system.  The class
is big, and its semiample contraction cannot be divisorial
because it lies in the interior of the movable cone.  Indeed
a nef class vanishing on curves covering a contracted divisor
lies on the boundary of the movable cone, as follows by
intersecting those curves with effective divisors avoiding
that divisor. Its ample-model map $N^{\mathrm c}\to Z^{\mathrm c}$
is therefore small.
The log-canonical divisor is crepant for this last contraction,
so the discrepancy bound passes to $Z^{\mathrm c}$. Its small
$\Q$-factorializations have precisely the common prime
valuations of the link. In particular, $N^{\mathrm c}\to Z^{\mathrm c}$
is one such small $\Q$-factorialization, and
$N^{\mathrm c}\dashrightarrow N_i$ is small for $i=1,2$.
Composing with the contraction towers
$N_i\to V_i\to B_i\to Q$ in the wall diagram gives exactly
the rational contractions asserted in the statement.

When $\Phi=0$, exceptional valuations above $W$ already have
log discrepancy greater than one.  A divisor of $W$ that is
contracted has a strictly negative coefficient in
\eqref{bd:ample-comparison}: intersect with a general curve
in its positive-dimensional contracted fibre, on which $A$
has positive degree.  Its coefficient in
\eqref{bd:original-comparison} is consequently strictly
positive.  This gives terminality also for these remaining
exceptional valuations.  The same reasoning applies to each
small model as above.
\end{proof}

To compare higher smooth starts, we use the following SNC discrepancy
estimate. Let $(R,\Gamma)$ be a smooth pair with SNC support and rational
coefficients $b_i<1$, allowing negative coefficients, and let $F$ be an
exceptional divisorial valuation. At the generic point of its center of
codimension $r$, choose regular parameters $x_1,\ldots,x_r$ containing
the equations of the boundary components through that center. Then
\begin{equation}\label{bd:snc-discrepancy}
 a(F,R,\Gamma)\geq
 \sum_{i=1}^r(1-b_i)\operatorname{ord}_F(x_i),
\end{equation}
where $b_i=0$ for parameters not defining a boundary component.
This follows from the Jacobian formula for a divisorial valuation
over a regular local ring, followed by subtracting the boundary
orders. All orders are positive integers.

\begin{lemma}[Higher common starts preserve the original marks]
\label{bd:higher-start}
Let $(W,\Phi)$ be smooth projective with SNC boundary whose
coefficients belong to $\{0,c\}$, where $0<c<1/4$ and
$\dim W\leq3$.  Suppose two Mori fibre spaces are outputs of
negative $(K_W+\Phi)$ programs. On a higher smooth common
start resolving the maps to both endpoints, keep the marks on
original prime divisors and assign
coefficient either $0$ or $c$ to any new prime divisor, with
SNC total support.  The same two outputs remain outcomes of
the program for this new marked pair.  In its geography the
original marked discrepancy bound can be taken to be
$\epsilon=1-c$, independently of the higher start.
The analogous assertion holds for ordinary terminal Mori
fibre spaces with zero boundary and $\epsilon=1$, including
the four-dimensional paths used below.
\end{lemma}

\begin{proof}
Write $(W',\Phi')$ for the higher smooth start with its new marking.
On the original smooth SNC pair, every exceptional valuation
has log discrepancy greater than one: \eqref{bd:snc-discrepancy}
gives the bound $2(1-c)>1$ for rational $c$. For irrational $c$,
apply it to rational coefficients decreasing to $c$ and use the
affine dependence of log discrepancy on the boundary coefficients.
The non-exceptional divisors
have discrepancy at least $1-c$.  These facts also hold on
the new smooth SNC start directly, without comparing the
number of its exceptional divisors.

Fix one endpoint $V$ and write $f:W'\to V$ for the resolved map.
For a divisor of $W$ contracted by
its negative program, discrepancy strictly improves.  If
its retained marking is $b\in\{0,c\}$, then on the higher
start its coefficient in the difference from the pullback
of $K_V+\Phi_V$ is
\[
                         a(F,V,\Phi_V)-1+b>0.
\]
For a divisor exceptional over $W$, monotonicity gives
$a(F,V,\Phi_V)\geq a(F,W,\Phi)>1$, and adding a new mark
$b\geq0$ again makes this coefficient positive.  Thus the
entire difference over $V$ is effective with strictly
positive coefficients on every endpoint-exceptional divisor.

Run the relative program over $V$.  Existence follows from
relative bigness for a birational morphism
\cite[Theorem~1.2]{BCHM10}.  At a relative minimal model the
exceptional difference is nef over $V$ and hence nonpositive
by negativity.  Since all the above positive exceptional
coefficients must have been contracted, the resulting map
to $V$ is small.  A projective small birational morphism to
a $\Q$-factorial variety is an isomorphism.  This realizes
the same endpoint, after which its Mori contraction can be
used.  The argument works for both endpoints.

We spell out the endpoint perturbations from
\cite[Lemma~4.1]{HM13}. Write $f_i:W'\to V_i$ for the two resolved
maps and $\pi_i:V_i\to B_i$ for their Mori contractions.
Choose small general effective ample rational
divisors $A,H_1,\ldots,H_s$, with the classes of the $H_j$ spanning
$N^1(W')$, and put $H=A+\sum_j H_j$. If $c$ is irrational, also
choose a rational boundary $\Phi_0\leq\Phi'$ sufficiently close to
$\Phi'$; for rational $c$ take $\Phi_0=\Phi'$. Make these choices so
that $A+\Phi_0-\Phi'$ is ample, both $f_i$ are
$(K_{W'}+\Phi_0+H)$-negative, and
$-(K_{V_i}+(f_i)_*\Phi_0+(f_i)_*H)$ is $\pi_i$-ample.
The exceptional differences have finitely many strictly positive
coefficients, which stay positive under these small changes;
the ampleness conditions are open as well.
For a sufficiently positive ample rational
divisor $C_i$ on $B_i$, the divisor
\[
 L_i=-(K_{V_i}+(f_i)_*\Phi_0+(f_i)_*H)+\pi_i^*C_i
\]
is ample. Choose a general effective rational representative
$F_i^{V_i}\sim_{\Q}L_i$ and put $F_i=f_i^*F_i^{V_i}$. Then
\[
 \Theta_i=\Phi_0+H+F_i,\qquad
 K_{V_i}+(f_i)_*\Theta_i\sim_{\Q}\pi_i^*C_i.
\]
Here $\Theta_i-\Phi'$ is ample: it is the sum of the ample divisor
$A+\Phi_0-\Phi'$ and the nef divisors $H_j,F_i$. Simultaneous general
representatives of sufficiently divisible multiples make the
dominating pair $(W',\Phi'+H+F_1+F_2)$ klt.
The negativity of $f_i$ therefore makes $B_i$ the required endpoint
ample model. The spanning classes $H_j$ and general two-dimensional
slice in \cite[Lemma~4.1]{HM13} complete these two perturbations to
one geography, retaining $\Phi'$ as the fixed original boundary.
Lemma~\ref{bd:central} now uses the unchanged lower bound
$1-c$.  With zero boundary the same proof starts from the
strictly positive discrepancies of terminal endpoints.
\end{proof}

\subsection{Bounded marked pairs and their contraction diagrams}

The central models preserve the discrepancy bound independently of the
chosen start. To bound the resulting links, we must control more than
these varieties: the contraction maps and their marked bad loci determine
the branch locations of the finite covers used later. We now prove this
upgrade for bounded log-Fano pairs. Theorem~\ref{bd:diagrams} will combine
it with Birkar's boundedness theorem for the central models.

In a \emph{marked contraction diagram}, the marks are finitely many
closed subsets, including specified divisors, and the arrows are
birational contractions or contraction morphisms. Diagrams
are considered up to compatible isomorphisms of their vertices; the
birational arrows use the common identification of function fields.
More precisely, every vertex field is embedded in the field of the
initial variety and every arrow respects these embeddings; in
particular, a cycle of birational arrows induces the identity on
that field.
Boundedness requires families for the vertices, marks, and graphs of
the arrows. In the contraction diagrams constructed below, the arrows
are the ample-model maps with these field identifications.
For a rational contraction we record both its factorization arrows
and the graph of their composite.
Graph resolutions will be chosen from the resulting
finite-type families. Repetition of a vertex or arrow does not add new data.

We first address the family issue that occurs when boundedness is
initially given without a relative $\Q$-Cartier divisor.

\begin{lemma}[Stratification of the marked log-Fano data]
\label{bd:cartier-strata}
Let $\mathcal P$ be a bounded collection of marked projective klt
log-Fano pairs $(Y,B)$, with coefficients of $B$ in a fixed finite
subset of $\Q$.  Finitely many stratified families, after finite
base changes of their parameter spaces, can be chosen so that:
\begin{enumerate}[label=\textup{(\roman*)}]
\item the log-canonical divisor is relatively $\Q$-Cartier, with a
common Cartier multiple on each stratum;
\item there is a simultaneous log resolution of the marked support,
and the coefficients of the crepant pullback of $B$ are constant on
each stratum;
\item any further specified rational Weil divisor, with bounded
marked support and coefficients in a fixed finite subset of $\Q$,
which is $\Q$-Cartier for every member has a common Cartier multiple
on each stratum, and its exceptional pullback coefficients are constant.
\end{enumerate}
The same assertions about $\Q$-Cartier classes hold for a bounded
collection of varieties with rational singularities whose smooth
projective resolutions have $H^1(\mathcal O)=H^2(\mathcal O)=0$.
In these families $\Q$-factoriality is a constructible condition.
\end{lemma}

\begin{proof}
We give the generic argument, and then apply noetherian induction.
Take the closure of the parameter points under consideration and an
irreducible component of this closure.  After shrinking and a
generically finite base change, we have flat normal geometric fibres
$\mathcal Y_t$, a family $\pi:\mathcal R\to\mathcal Y$ of resolutions,
and smooth projective fibres $\mathcal R_t$ on which the marked strict
transforms and exceptional divisors form a relative SNC divisor.
The individual exceptional components can also be separated by this
base change.  These are the usual spreading operations on a resolution
of the geometric generic fibre.

For the members in question, rational singularities and vanishing
give
\[
 H^i(\mathcal Y_t,\mathcal O_{\mathcal Y_t})
 =H^i(\mathcal R_t,\mathcal O_{\mathcal R_t})=0\quad(i=1,2).
\]
For log-Fano pairs, the downstairs vanishing follows from
Kawamata--Viehweg vanishing and the equality follows from rational
singularities; see
\cite[Theorem~3.13.1 and Corollary~5.7.7]{FujinoFoundationsDraft057}.
The parameter points of these members are Zariski dense in the
closure chosen above. Thus they satisfy the dense-subset hypothesis
of the relative Picard-space theorem \cite[Theorem~1.3]{CLZ26}:
rational singularities and $H^1(\mathcal O)=H^2(\mathcal O)=0$
are required on that dense subset. Upper semicontinuity also gives
the displayed vanishings on an open set. The theorem permits a further generically finite
base change and shrinking such that the rational numerical divisor
spaces of $\mathcal Y_t$ and $\mathcal R_t$, their pullback map, and
the exceptional divisor classes are identified with fixed rational
linear data.  In particular, choose relative Cartier divisors whose
classes span the downstairs space on every fibre.

Here is the required $\Q$-Cartier criterion.  For a rational Weil
divisor $D$ on a fibre, let $\widetilde D$ be its strict transform.
Then
\begin{equation}\label{bd:cartier-test}
 D\text{ is }\Q\text{-Cartier}
 \quad\Longleftrightarrow\quad
 [\widetilde D]\in
 \pi^*N^1(\mathcal Y_t)_{\Q}
       +\sum_E\Q[E]\ \subset N^1(\mathcal R_t)_{\Q}.
\end{equation}
The forward implication follows by pullback.  For the converse,
numerical equivalence of rational divisors on $\mathcal R_t$ is
rational linear equivalence: its $\operatorname{Pic}^0$ vanishes,
and the numerically trivial part of the N\'eron--Severi group is
torsion.  Thus the asserted relation, after clearing denominators,
is a linear-equivalence relation with a Cartier divisor pulled back
from $\mathcal Y_t$ and exceptional divisors.  Pushing down proves
that a multiple of $D$ is linearly equivalent to a Cartier divisor.
This is precisely the converse in \eqref{bd:cartier-test}.

Apply the criterion to the fixed rational classes on the resolution.
For the log-canonical class one may use
$K_{\mathcal R_t}+\widetilde B_t$ modulo exceptional classes; this
does not presuppose that $K_{\mathcal Y_t}+B_t$ is $\Q$-Cartier.
The membership in \eqref{bd:cartier-test} holds on the dense set of
members under consideration, and the linear data are fixed.  It
therefore holds at the geometric generic point, with fixed rational
coefficients in a basis of relative Cartier classes.  Rational linear
equivalence upstairs, followed by pushdown, proves that the generic
log-canonical divisor is $\Q$-Cartier.  Clear denominators and spread
this linear equivalence.  Its possible vertical errors disappear
after shrinking, giving a relative Cartier multiple.  This argument
applies simultaneously to each additional specified class.

Exceptional coefficients in a pullback are unique.  Indeed, an
exceptional rational divisor numerically trivial over its target is
zero, by applying the negativity lemma to it and its negative.
They are thus the solutions of fixed rational linear equations on
the resolution, and are constant after shrinking.  The klt condition
is now the finite set of strict inequalities saying that the crepant
SNC coefficients are less than one.  Relative log-anticanonical
ampleness is open.  Thus an open set consists entirely of the required
log-Fano pairs, with all the asserted properties.

Finally, all Weil divisors downstairs are $\Q$-Cartier exactly when
the right side of \eqref{bd:cartier-test} exhausts
$N^1(\mathcal R_t)_{\Q}$: every divisor upstairs pushes to a Weil
divisor, and the same argument applies.  This is again a condition
on the fixed linear data.  Applying the construction to the
remaining closed parameter sets proves all the assertions by
noetherian induction.
\end{proof}

Choi--Li--Zhou prove boundedness of the birational contraction models
of members of a fixed bounded family of Fano-type varieties
\cite[Theorem~1.7 and Section~5.1]{CLZ26}.
We use their fibrewise small $\Q$-factorializations and the comparison
between relative and fibre MMPs
\cite[Theorems~1.4--1.6 and Lemma~5.5]{CLZ26}.
The additional data in the following statement are the marked supports,
graphs, exceptional loci and their resolutions.
The proof uses finite generation to list the canonical contraction
maps on a generic member, the cited family results to retain the
complete list after shrinking, and then spreading to include their
graphs and marks.

\begin{lemma}[Uniform contraction diagrams]
\label{bd:contractions}
For a bounded collection of projective klt log-Fano pairs with fixed
finite coefficient set, all diagrams obtained from a small
$\Q$-factorialization by small modifications and rational contractions
are bounded.  The assertion includes the transforms and images of a
bounded marked support, the exceptional loci of the arrows, and their
chosen graph resolutions.  It also includes a non-extracting birational map
followed by a contraction morphism.
Here a diagram is a subdiagram of the collection of contraction models
and their canonical arrows compatible with the fixed function field.
\end{lemma}

\begin{proof}
Lemma~\ref{bd:cartier-strata} supplies families of Fano-type fibres.
Take smooth parameter strata and remove the images of any vertical
exceptional components of the simultaneous resolution.  The relative
$\Q$-Cartier equality and the crepant coefficients less than one
then make the total-space pair klt as well.  This verifies the
total-space hypothesis of \cite[Theorem~1.4]{CLZ26}.
After a generically finite base change and shrinking,
\cite[Theorems~1.4--1.6]{CLZ26} gives fibrewise small
$\Q$-factorializations and compatible numerical spaces, nef cones,
Mori chambers, and MMP steps.  Its hypotheses hold: the fibres have
rational singularities and the stated vanishings, and every fibre
is of Fano type.  Thus the very-general-fibre alternative in
\cite[Theorem~1.2]{CLZ26} applies.

For clarity, the finiteness-to-boundedness step is as follows.  On
the geometric generic fibre of one stratum choose a small
$\Q$-factorialization.  Finite generation supplies finitely many
small models, and finitely many faces and semiample contractions
on each of them.  All this finite data spreads after a finite
extension and shrinking.  The two inclusions needed for constancy
of a nef cone can be seen directly: spread its finitely many
semiample generators, which remain semiample on an open set, and
spread effective curves generating the dual polyhedral cone.
The divisor generators give one inclusion; the curve generators
give the reverse.  The fibrewise $\Q$-factorializations and
compatible numerical spaces ensure that these comparisons take
place in the same space.  The relative chamber/MMP compatibility
cited above then rules out further chambers on a fibre.
Here the source has first been made fibrewise $\Q$-factorial, so the
partial-MMP chamber description agrees with that from small modifications;
the converse direction in \cite[Theorem~1.6(3)]{CLZ26} applies.  The
generic finite list of graphs consequently supplies the complete
list on an open set.  Apply noetherian induction to the remaining
strata.  This proves boundedness of the diagrams, not just of their
vertices.

We also check the last assertion.  If $f:Y\dashrightarrow Z$ is a
non-extracting birational map, take a sufficiently ample Cartier
divisor $H$ on $Z$ and a general member avoiding the generic points
of the finitely many centers of divisors contracted by $f$.
Write $D$ for its strict transform on $Y$. The corresponding linear
system has no fixed prime component.
Moreover,
\[
                 R(Y,D)=R(Z,H).
\]
At a prime divisor retained by the map this is the same valuation
test on both sides.  At a prime divisor contracted on $Y$, a
regular section of $\mathcal O_Z(mH)$ pulls back regularly, and
the chosen $H$ has order zero at that valuation.  Normality then
gives equality of the section spaces.  Hence $Z$ is the ample
model of a movable system on $Y$.  A subsequent contraction is
handled by taking the pullback of an ample divisor from its base.
The Mori dream space chamber description therefore includes these
maps too.
If an intermediate target is not $\Q$-factorial, apply this description
to the composite from the initial $\Q$-factorial source; the diagram
also records the intermediate maps.

Finally take closures of the marked transforms on the finitely
many graph families, and resolve their union with the exceptional
loci.  Proper images, inverse images, and irreducible components
of these finite-type families remain bounded after stratification.
This gives the stated marked and exceptional data.
\end{proof}

\begin{remark}[Bad loci in a bounded diagram]\label{bd:bad-loci}
For any of these finitely many contraction families, the proper
closed locus outside which a specified generic fibre property holds
can be chosen in the same bounded diagram.  For example, use the
smooth locus in a conic or del Pezzo family.  More generally, spread
a resolution of the geometric generic fibre and its very free
curves when that fibre is rationally connected.  Smoothness and
deformation of the curves give a dense good open.  Shrinking and
stratifying ensures that it is dense in each relevant fibre of the
parameter family.  The complementary closed sets and their
preimages are then bounded.  A prime divisor contained in such a
proper closed preimage is an irreducible component of that
preimage, so this observation bounds the divisors actually used
below; it makes no assertion about arbitrary subvarieties inside
a bounded variety.
\end{remark}

\subsection{Uniform geometric generic link diagrams}

\begin{theorem}[Uniform geometric generic link diagrams]
\label{bd:diagrams}
Bounded marked families of geometric generic diagrams can be chosen
with the following properties:
\begin{enumerate}[label=\textup{(\roman*)}]
\item There are such families depending only on the dimension so
that, for every specified birational map
$f:V/B\dashrightarrow V'/B'$ between terminal four-dimensional
Mori fibre spaces, there is an ordinary Sarkisov factorization of
$f$ whose individual geometric generic link diagrams belong to
these families.
\item For every fixed rational $0<c<1/4$, there are such families
depending only on the dimension and $c$ so that, for every smooth
projective threefold $W$ with a reduced SNC divisor $D$ and every
two given negative $(K_W+cD)$ programs with Mori fibre-space outputs,
there is a connecting log Sarkisov path whose individual geometric
generic link diagrams belong to these families.
\end{enumerate}
In \textup{(i)} the path induces the specified map $f$; in
\textup{(ii)} it induces the birational map determined by the two
given maps from $W$. The diagrams retain these identifications of
their function fields.
For a link with common base $Q$, put $F=\C(Q)$ and choose an
algebraic closure $\overline F$. The bounded diagrams are the
restrictions to the geometric generic fibre over $\overline F$.
They include the small models, both ends,
their base contractions, exceptional loci, marked transforms,
and bounded bad loci as in Remark~\ref{bd:bad-loci}.
In \textup{(i)} the marks are the exceptional and bad loci of the
link contractions; in \textup{(ii)} they also include the restrictions
of the transforms of the original divisor $D$.
Each induced end-base contraction $B\to Q$ is of Fano type
over $Q$.  Its geometric generic fibre is integral and has
a rationally connected smooth proper model.
In \textup{(ii)} all marked prime divisors on a link model are
transforms of primes on the chosen common start: the maps
from that start do not extract.  Higher starts with the mark
rule of Lemma~\ref{bd:higher-start} are allowed.

In both parts the families are independent of the start, its
resolution, the birational map,
the length of a path, or the K3 test collection.
\end{theorem}

\begin{proof}
Apply the geography above to a common smooth resolution of the
specified map in \textup{(i)}, and to the given common-start maps in
\textup{(ii)}, retaining their function-field identifications. Lemmas
\ref{bd:central} and \ref{bd:higher-start} give a central ample
model $Z^{\mathrm c}/Q$, reached by a small map from
$N^{\mathrm c}/Q$, with $\epsilon=1$ in the first case
and $\epsilon=1-c$ in the second. Put
$Z=Z^{\mathrm c}_{\overline F}$. Its dimension is at most four,
the negative log-canonical divisor is ample, and the same
discrepancy bound holds.  Birkar's boundedness theorem
\cite[Theorem~1.1]{BirkarBAB21} applies: projective varieties
admitting an $\epsilon$-lc boundary with nef and big negative
log-canonical divisor form a bounded family. In particular
it bounds the underlying varieties $Z$ just constructed.

The marked boundary is bounded as well. In case \textup{(ii)},
write its restriction to $Z$ as $cD_Z$, and let $H$ be one of
the bounded very ample polarizations supplied by boundedness.
In fibre dimension $r$,
\[
              cD_Z\cdot H^{r-1}\leq -K_Z\cdot H^{r-1}.
\]
The right side is bounded by the coefficient of the Hilbert
polynomial of the bounded polarized varieties. Thus $D_Z$
has bounded degree.  Reduced cycles of bounded degree in a
fixed projective space form a bounded family, so the marked
central pairs are bounded. Marks vertical over $Q$ have empty
restriction to the generic fibre. Horizontal components may split
over $\overline F$, but their nonzero coefficients remain $c$
and their total degree is still bounded.

All the geometric generic link vertices and bases are
rational contractions of small $\Q$-factorializations of
these central pairs.  Lemma~\ref{bd:contractions} supplies
bounded diagrams, and Remark~\ref{bd:bad-loci} supplies the
closed bad loci.  Small loci do not acquire divisorial
components under the algebraic extension to the geometric
generic field. The same constructions apply over algebraically closed
fields of characteristic zero. For the family results stated over
$\C$, we use their finite-data spreading arguments and algebraically
closed base change to obtain the corresponding geometric generic
families here.
Every bound was taken from one of these fixed bounded
families, not by accumulating constants along a path.

For the assertion about base contractions, the central rank-two
model is of Fano type over $Q$ by Lemma~\ref{bd:central}.
This property survives small modifications and descends under
contractions, including fibre-type contractions
\cite[Lemma~2.8, with the relative convention after
Lemma--Definition~2.6]{ProkhorovShokurov09}.
It therefore holds for every end and its base $B\to Q$.
These are contractions of normal varieties, so their generic
fibres are geometrically integral in characteristic zero.
The rational connectedness theorem for klt log-Fano
contractions \cite{HM07Vertical}, applied also to a resolution,
gives the asserted rationally connected smooth proper models.
\end{proof}

\subsection{The low-discrepancy extractions}

The link diagrams have now been bounded. A different issue arises in
Proposition~\ref{sf:conic-end}: a conic base is controlled on its geometric
generic fibre, but an arbitrary smooth resolution need not be bounded
even there. The relative log MMP over that base will leave a geometric
generic model extracting only valuations of log discrepancy at most one,
as proved in that proposition. We prepare the bound for that surviving
model here; we do not bound the resolution itself.

The contraction lemma applies to maps which do not extract divisors.
We therefore first extract all exceptional prime valuations of log
discrepancy at most one. Any model extracting only some of them is a
birational contraction of this common model and falls under that lemma.

\begin{lemma}[Finite extraction list on an SNC pair]
\label{bd:snc-places}
Let $(R,\Gamma)$ be a smooth pair with SNC support and rational
coefficients $b_i<1$, allowing negative coefficients.  There are
only finitely many exceptional divisorial valuations $F$ with
$a(F,R,\Gamma)\leq1$.  They are obtained by a finite collection
of blow-ups of boundary strata.  If the pairs and their SNC
components vary in a fixed log-smooth family with a fixed finite
coefficient set, all these valuations occur in finitely many such
relative constructions after stratification.
\end{lemma}

\begin{proof}
At the generic point of the center of $F$, of codimension $r$, use the regular parameters
and boundary coefficients in \eqref{bd:snc-discrepancy}.
If the center is not
a boundary stratum, at least one $b_i$ is zero; since an
exceptional center has $r\geq2$, the right side is strictly
greater than one.  Thus a valuation in question centers at a
boundary stratum.  Write its orders there as $w_i$.  They satisfy
\[
                     \sum_i(1-b_i)w_i\leq1.
\]
There are finitely many possibilities because every $1-b_i$ is
positive.

Blow up the stratum.  In a chart containing the center of the
valuation, choose an index with minimal order; the new orders
are that minimum and the differences from it.  Zero differences
correspond to units unless the center lies in a further proper
subvariety of the exceptional stratum.  The latter would be a
non-stratum center and is excluded by
\eqref{bd:snc-discrepancy} applied to the crepant SNC transform.
This transform is still klt, so its coefficients remain strictly less
than one and the same non-stratum exclusion applies at every step.
Thus, until the valuation itself becomes a divisor, its center
continues to be a stratum and the sum of its positive orders
strictly decreases.  This is the usual monomial subtraction
algorithm, and proves both that the valuation is a stratum
valuation and that finitely many blow-up sequences suffice.
In a fixed family there are finitely many strata and bounded
orders; separate their geometric components and perform the
same finite sequences relatively.
\end{proof}

\begin{theorem}[Bounded terminalization diagrams]
\label{bd:terminalization}
Let $(T,\Lambda)$ range over a bounded collection of projective
klt log-Fano pairs with fixed finite rational coefficient set and
bounded marked support.  There is a bounded collection of marked
diagrams containing a $\Q$-factorial crepant terminalization
extracting exactly the exceptional valuations
\[
                         a(F,T,\Lambda)\leq1.
\]
Every normal projective model which extracts only a subset of
these valuations, followed by non-extracting birational moves
and contraction morphisms, occurs among bounded marked
contraction diagrams of these terminalizations.  The bound
depends on the given family, not on a resolution chosen to
construct one of the models.
\end{theorem}

\begin{proof}
Use Lemma~\ref{bd:cartier-strata} to obtain simultaneous crepant
log resolutions.  Their coefficients are in a fixed finite set
strictly below one on each stratum.  The exceptional divisors
already on the resolution are a finite list; all additional
valuations of log discrepancy at most one are in the uniformly
finite list of Lemma~\ref{bd:snc-places}.  Extract all of them on
a common resolution.

The extraction theorem \cite[Corollary~1.4.3]{BCHM10} produces a
$\Q$-factorial model $\tau:T^{\mathrm t}\to T$ extracting exactly
this set.  The extra restriction concerning non-klt centers in
that theorem is empty here.  Its crepant boundary
$\Lambda^{\mathrm t}$ has coefficient $1-a(F,T,\Lambda)$ on an
extracted divisor, and hence is effective.  By construction the
pair is terminal.  Also
\[
 -(K_{T^{\mathrm t}}+\Lambda^{\mathrm t})
             =\tau^*\bigl(-(K_T+\Lambda)\bigr)
\]
is nef and big, so $T^{\mathrm t}$ is of Fano type.  Apply this
construction on a geometric generic member and spread it; the
fixed extraction list, crepant equality, and the discrepancy
inequalities persist after shrinking.  Thus the terminalizations
and all their marks form bounded families.  This construction
does not resolve an unbounded set of arbitrarily chosen higher
models.

For a model $T'$ extracting a subset, every prime divisor on
$T'$ is already a valuation appearing on $T^{\mathrm t}$.
Consequently $T^{\mathrm t}\dashrightarrow T'$ is a birational
contraction.  Lemma~\ref{bd:contractions}, applied to the bounded
Fano-type terminalization families, supplies its diagram and
every indicated subsequent contraction.  In applying that lemma
one can either use the displayed weak log-Fano pair and a small
ample perturbation of its boundary on each stratum, or spread
a Fano-type boundary chosen on its generic member.  Both give
families of klt log-Fano pairs on an open set.  Include all the
extracted divisors as marks, even those with crepant coefficient
zero.  The marked support on every output is then among the
bounded transforms already considered.
\end{proof}

\subsection{Choosing the ramification coefficient once}

\begin{lemma}[A uniform coefficient on bounded bases]
\label{bd:coefficient}
Let $(T,D)$ range over a bounded collection with $T$ a
projective klt Fano variety and $D$ a reduced divisor.  Assume
$K_T$ and $D$ are $\Q$-Cartier and $D$ is numerically
proportional to $-K_T$.  There is a rational $c_*>0$ such
that $(T,cD)$ is klt log Fano for every rational
$0<c<c_*$.  This also holds when the numerical proportionality
is known over a characteristic-zero field before geometric
base change and both divisors are defined over that field.
\end{lemma}

\begin{proof}
Zero-dimensional members have $D=0$ and impose no restriction on $c$;
we treat the positive-dimensional members below.
Apply Lemma~\ref{bd:cartier-strata} to the log-Fano pairs
$(T,0)$ with additional marked divisor $D$.
On finitely many simultaneous log resolutions the crepant
coefficients for $(T,0)$ and the coefficients of $\pi^*D$
are fixed.  Write $a_E=a(E,T,0)>0$ and $m_E$ for the
coefficient of $\pi^*D$ on an exceptional component.
Choosing $c<1/2$ and $cm_E<a_E/2$ for every $m_E>0$
makes every coefficient of the crepant pullback of $cD$
strictly less than one.  These finitely many inequalities
give a uniform positive klt bound for $c$.

For ampleness use a bounded very ample divisor $H$ and
let $r=\dim T$.  A common multiple $mK_T$ is Cartier
on each of the finitely many strata.  Thus
\[
 k=-K_T\cdot H^{r-1}>0,\qquad mk\in\Z_{>0}.
\]
The degree $e=D\cdot H^{r-1}$ is a bounded nonnegative
integer.  Numerical proportionality says
$D\equiv(e/k)(-K_T)$.  Since $e/k\leq me$, choosing
$cme<1/2$ for all these data gives
\[
 -(K_T+cD)\equiv (1-ce/k)(-K_T),
\]
which is ample.  Taking the minimum of the finitely many
rational choices gives $c_*$.

Finally numerical equivalence between divisors defined
over the original field persists geometrically.  Any
geometric curve is defined over a finite extension; all
its conjugates have the same intersection with such a
divisor.  Its orbit descends to a curve cycle over the
original field, so a zero intersection downstairs
forces zero on every conjugate.  The preceding argument
therefore applies after the geometric base change.
\end{proof}

For the next corollary, let $V\to T$ be a conic Mori fibre space and
let $\alpha\in\Br(\C(T))[2]$ be the class of its generic conic.
At a prime divisor $P\subset T$, its Brauer residue belongs to
$H^1(\C(P),\Z/2)$. Define $D_T$ to be the reduced sum of the primes
with nonzero residue; this is the \emph{full visible divisorial
ramification support} on this model. For a map $T\to Q$, put
$F=\C(Q)$. On $T_{\overline F}$ we retain the divisor
$(D_T)_{\overline F}$ obtained by restricting and extending this
original support. A residue can become zero after extension of
constants; its original supporting divisor remains marked.

\begin{corollary}[Order of the uniform choices]
\label{bd:choice-order}
For the conic-base ends $V/T$ of ordinary fourfold links in the
paths of Theorem~\ref{bd:diagrams}, over $Q$ with $\dim Q\leq1$,
a single rational number
\[
                         0<c<1/4
\]
can be fixed such that the geometric generic pair
$(T_{\overline F},c(D_T)_{\overline F})$, with $F=\C(Q)$, is klt
log Fano. Here $D_T$ is the full visible divisorial ramification
support defined on the original conic base above.
Fix this number before applying the log part of
Theorem~\ref{bd:diagrams} or
Theorem~\ref{bd:terminalization}.  All resulting
large-degree exclusions have a common threshold
independent of the K3 test subcollection and of the chosen
Sarkisov path.
\end{corollary}

\begin{proof}
First use only the ordinary part of
Theorem~\ref{bd:diagrams}. At the indicated end the tower is
$N_i\xrightarrow{e_i}V\to T\to Q$. Since $\dim T=3$ and
$\dim Q\leq1$, the last contraction has positive relative
dimension and $\rho(T/Q)\geq1$. Equation~\eqref{bd:rank-tower} gives
\[
 2=\rho(N_i/V)+1+\rho(T/Q).
\]
Consequently $\rho(N_i/V)=0$ and $\rho(T/Q)=1$: the arrow $e_i$
is the identity and $T\to Q$ is an elementary contraction of fibre
type. This calculation takes place over $Q$ before geometric
generic base change.
The geometric generic model $T_{\overline F}$ and its conic
projection occur in the bounded diagram. The restriction of the
original visible ramification is contained in the
bounded bad locus of this projection, and its divisorial
support is a union of the finitely many divisorial
components there.  These marked supports are therefore
bounded.

Apply \cite[Lemma~4.6]{Druel16} to the elementary conic Mori
contraction $V\to T$. Its source is terminal and $\Q$-factorial,
so its base $T$ is $\Q$-factorial.
Divisors on $T_F$ extend by closure to Weil divisors on $T$,
which are $\Q$-Cartier. Restriction therefore surjects onto the
generic numerical divisor space. The relative rank-one condition
and the nonzero restriction of an ample class give $\rho(T_F/F)=1$.
Fano type descends under contractions
\cite[Lemma~2.8]{ProkhorovShokurov09}; hence $T_F$ and
$T_{\overline F}$ are of Fano type. Over the original generic field
$F$, rank one implies that $-K_{T_F}$ is ample: it is the sum of
an ample log-anticanonical class and an effective class.
It also implies numerical proportionality of $(D_T)_F$ and
$-K_{T_F}$. These two divisors are $\Q$-Cartier before, and
hence after, base change. Their numerical proportionality persists
over $\overline F$ by Lemma~\ref{bd:coefficient}, even if the
geometric Picard rank increases. That lemma
provides a uniform $c_*$ from these ordinary diagrams.

Choose a rational $c<\min\{1/4,c_*\}$ once.  With this
fixed value the marked log links have the uniform bound
of Theorem~\ref{bd:diagrams}, and the bounded
$(T_{\overline F},c(D_T)_{\overline F})$ data have the extraction bound of
Theorem~\ref{bd:terminalization}.  Only finitely
many bounded families, dimensions, and bounded-degree
cover constructions are used subsequently: the exceptional and
bad-locus components in Section~\ref{sl:section}, and the
ramification and del Pezzo class covers in Section~\ref{sf:section}.
Their branch and degree bounds are checked there. Take the
maximum of their thresholds in
Theorem~\ref{bd:exclusions}.  Every bound depends
only on those families and this chosen $c$, proving
the asserted independence.
\end{proof}
